\documentclass[a4paper,12pt]{article}

\usepackage[
    left=1.8cm,
    right=1.8cm,
    top=1.5cm,
    bottom=1.5cm
]{geometry}
\usepackage[T1]{fontenc}
\usepackage[utf8]{inputenc}
\usepackage[french]{babel}
\usepackage{array}
\usepackage{graphicx}
\usepackage{xcolor}

\pagestyle{empty}


\begin{document}

\begin{center}
    \LARGE{Soutenance de projet} \\
    \large{LU2IN013 --- 6 Mai 2026}
\end{center}


\section*{Consignes}

En groupe.

\vspace{0.2cm}

Slides disponibles sur le repo GitHub \textbf{la veille à minuit} (sinon -2 pts).

\vspace{0.2cm}

\textbf{Durée de présentation :} 12 min (dont 5 min de démo) + 8 min de questions

\vspace{0.2cm}

\textbf{Temps de parole égaux. On ne lit pas ses notes.}

\vspace{0.2cm}

Contenu:
\vspace{0.2cm}
\begin{itemize}
\item[$\ast$] Replacer le projet en 1 min max dans un contexte général:
développements récents de l'IA et de la robotique.
Penser aux applications potentielles d'un tel robot.
\item[$ $]
\item[$\ast$] Schéma fonctionnel avec les différents modules/composants
(simulation, contrôleur, stratégies, robot réel, robot simulé, obstacles simulés, etc),
et montrant comment ils intéragissent entre eux.
\item[$ $]
\item[$\ast$] Schéma de l'architecture logicielle avec les différentes classes
et les principales méthodes associées à ces classes.
Un diagramme UML simplifié peut faire l'affaire.
\item[$ $]
\item[$\ast$] Réflexion sur l'ajout d'une ou deux fonctionnalités intéressantes.\\
Répondre à la question : comment les inclure dans l'architecture logicielle ? Sous forme de classes
ou de fonctions ? Avec quelles classes, fonctions ou méthodes existantes intéragiraient ces nouveaux 
composants ? Montrer comment l'on ferait évoluer le logiciel (capacité de planification à long terme).
\item[$ $]
\item[$\ast$] Retour d'expérience. Qu'a-t-on appris ?
\end{itemize}


\section*{Recommandations}

\begin{itemize}
\item[$\ast$] Numéroter les slides.
\item[$\ast$] Vocabulaire précis. Titres précis. Le vent est irritant.
\item[$\ast$] DES SCHÉMAS !!! Pas de bullet points pourris. 
Bullet points tolérés en conclusion, intro ou à côté d'un schéma ou d'une image.
\item[$\ast$] Pas trop de texte. Quand c'est trop dense, on se perd.
\item[$\ast$] SURTOUT PAS de captures d'écran de code
\item[$\ast$] Les slides doivent servir le discours.
\item[$\ast$] Votre présentation raconte une histoire et doit avoir un 
fil conducteur, lequel permet de structurer votre pensée et votre présentation.
\item[$\ast$] \textit{Simplicity is better.}
\item[$\ast$] Écrivez votre \textit{speech} à l'avance. 
MAIS vous n'aurez pas le droit de lire vos notes.
\end{itemize}

\end{document}
